\documentclass[a4paper,11pt]{article}

\usepackage[utf8]{inputenc}
\usepackage[T1]{fontenc}
\usepackage{geometry}
\usepackage{xcolor}
\usepackage{listings}

\geometry{margin=1in}

\lstset{
  basicstyle=\ttfamily\small,
  backgroundcolor=\color{gray!10},
  frame=single,
  breaklines=true,
  tabsize=2
}

\title{Maintenance and Troubleshooting}
\author{Group 2 -- Byte Me}
\date{}

\begin{document}

\maketitle

\section{Common issues}

\subsection{Backend won't start}

Usually a database problem. Check that Postgres is running and that the \texttt{DB\_URL}, \texttt{DB\_USERNAME}, and \texttt{DB\_PASSWORD} in your environment match your local setup. The default URL is \texttt{jdbc:postgresql://localhost:5432/byteMe}. If the database doesn't exist yet, create it with \texttt{createdb byteMe}.

If Flyway fails on startup, it's probably a migration conflict. You can wipe the schema and let it rebuild, or set \texttt{FLYWAY\_ENABLED=false} temporarily to skip migrations.

Port 8080 already in use? Either kill the other process or set \texttt{SERVER\_PORT} to something else.

\subsection{CORS errors in the browser}

The backend only accepts requests from origins listed in \texttt{CORS\_ORIGINS}. Locally this defaults to \texttt{http://localhost:3000}. If you're running the frontend on a different port or domain, update the variable. Watch out for trailing slashes and protocol mismatches, they matter.

\subsection{JWT / auth problems}

\texttt{JWT\_SECRET} needs to be at least 32 bytes long. If it's too short the backend will throw a \texttt{SignatureException} on startup or when generating tokens.

Expired tokens just result in unauthenticated requests, you won't get a specific error message. If a user suddenly can't access anything, have them log out and back in.

\subsection{Frontend won't build}

Run \texttt{npm install} first. If the build fails on TypeScript errors, check that you're on Node 20+. Older versions sometimes have issues with React 19 types.

\subsection{Tests failing}

Backend tests use an in-memory H2 database, so they don't need Postgres running. If they fail it's usually a code issue, not config.

Frontend tests need the dev dependencies installed. Run \texttt{npm install} and then \texttt{npx vitest run}. If you get module resolution errors, delete \texttt{node\_modules} and reinstall.

\section{Routine maintenance}

\subsection{Updating dependencies}

Backend: bump versions in \texttt{pom.xml} and run \texttt{mvn clean test} to check nothing breaks.

Frontend: run \texttt{npm update} or edit \texttt{package.json} directly, then \texttt{npm install} and \texttt{npm run build}.

\subsection{Adding a database migration}

Create a new file in \texttt{backend/src/main/resources/db/migration/} following the naming convention: \texttt{V3\_\_description.sql}. Flyway picks it up automatically on the next backend startup.

Don't edit existing versioned migrations. If you need to change something from V1 or V2, create a new migration that alters the table.

\subsection{Adding a new API endpoint}

Create or edit a controller in \texttt{backend/src/main/java/com/byteme/app/}. If the endpoint should be public (no login required), add its path to the \texttt{SecurityConfig} permit list. Otherwise it's protected by default.

\subsection{Redeploying}

Push to \texttt{master}. Railway picks it up and rebuilds both services. Check the Railway dashboard for build logs if something fails. You can roll back to a previous deployment from the dashboard.

\section{Logs}

Backend logs go to stdout. Locally you'll see them in the terminal where you ran \texttt{mvn spring-boot:run}. On Railway, check the service dashboard for real-time logs.

Log levels are configurable:

\begin{itemize}
  \item \texttt{LOG\_LEVEL\_SECURITY} -- Spring Security logs (default: WARN)
  \item \texttt{LOG\_LEVEL\_SQL} -- Hibernate SQL queries (default: WARN)
\end{itemize}

Set either to \texttt{DEBUG} if you need more detail.

\section{Contacts}

If something is broken in production, check the Railway dashboard first. The deployment history shows whether the last deploy succeeded. If the services are up but the site is misbehaving, check the backend logs for stack traces.

For code-level issues, the GitHub repo has the full commit history and issue tracker.

\end{document}
